\chapter{Matching independent labels on the binary tree}
\label{ch:iid}

This chapter proves the matching theorem for independent labels drawn from a countable graph.
The proof reduces matching on a binary tree to the iteration of a one-site potential, proves
that this potential contracts under the binary-tree recursion, and then passes from finite
heights to the infinite tree. The entire argument is formalised in \texttt{GraphMatching},
which depends on Mathlib alone.

\section{Set-up and the compatibility rule}

\begin{definition}[The finite binary tree and its automorphisms]
\label{def:aut}
\lean{GraphMatching.Aut}
\leanok
For $h\geq0$ the complete binary tree $\bbB_h$ consists of the words over $\set{1,2}$ of
length at most $h$, arranged in levels $0$ through $h$. Its $2^h$ leaves are the words of
$\set{1,2}^h$. An automorphism $\pi\in\Aut(\bbB_h)$ fixes the root and, at every internal
vertex, either fixes or interchanges the two child subtrees, acting recursively inside
them. In particular $\pi$ preserves levels.
\end{definition}

The induction on the height uses this recursive description, which is also the definition of the
group in Lean. The declaration \texttt{fullSim\_iff\_graphAut} verifies that these recursively
defined maps are precisely the root-fixing graph automorphisms of the adjacency graph of
$\bbB_h$.

\begin{definition}[Compatibility]
\label{def:compat}
\lean{GraphMatching.compat}
\leanok
Let $G=(V,E)$ be a countable undirected graph. Two labels $v,w\in V$ are \emph{compatible}
when $d_G(v,w)\leq1$, so that the labels compatible with $v$ are exactly those of the
closed unit ball $B_G(v,1)$. The relation is symmetric and reflexive, and those are the
only two properties the proof uses.
\end{definition}

\begin{definition}[Leaf and full matching]
\label{def:matching}
\lean{GraphMatching.leafSim, GraphMatching.fullSim}
\leanok
\uses{def:aut, def:compat}
A leaf labelling of $\bbB_h$ is a map $x\colon\set{1,2}^h\to V$ and a full labelling is a
map $x\colon\bbB_h\to V$. Two leaf labellings are matched when some $\pi\in\Aut(\bbB_h)$
carries $x$ to a labelling compatible with $y$ at every leaf, and likewise for full
labellings at every vertex.
\end{definition}

\begin{definition}[The weight and the one-site potential]
\label{def:potential}
\lean{GraphMatching.phiA, GraphMatching.PhiA}
\leanok
For $\alpha\geq1$ write
\[
 \phi_\alpha(t):=\frac{t}{(1-t)^\alpha}\qquad(0\leq t<1),
\]
so that $\phi_\alpha\geq0$ and $\phi_\alpha(t)\geq t$. Let $(\cX,\mu)$ be a countable
probability space carrying a symmetric reflexive relation $R$, and put
\[
 q(x):=\mu\set{y\colon x\mathrel{\not R}y},\qquad r(x):=1-q(x),
\]
with $S_\mu:=\set{x\in\cX\colon\mu\set{x}>0}$ the support. Reflexivity gives $r(x)\geq
\mu\set{x}>0$, hence $q(x)<1$ on $S_\mu$, and the \emph{one-site potential}
\[
 \Phi(R,\mu):=\bbE\,\phi_\alpha(q(X))
 =\sum_{x\in S_\mu}\mu\set{x}\,\phi_\alpha(q(x)),\qquad X\sim\mu,
\]
is a well-defined element of $[0,\infty]$.
\end{definition}

\begin{definition}[The symmetrised square]
\label{def:square}
\lean{GraphMatching.SquareRel}
\leanok
\uses{def:potential}
On $\cX^2$ with the product measure $\mu^2$, define $R^\square$ by
\[
 (x_0,x_1)\mathrel{R^\square}(y_0,y_1)
 \iff
 \bigl(x_0\mathrel{R}y_0\wedge x_1\mathrel{R}y_1\bigr)
 \vee
 \bigl(x_0\mathrel{R}y_1\wedge x_1\mathrel{R}y_0\bigr).
\]
It is again symmetric and reflexive, with support $S_\mu\times S_\mu$.
\end{definition}

\begin{definition}[The graph potential]
\label{def:etaG}
\lean{GraphMatching.etaGA}
\leanok
\uses{def:compat, def:potential}
Write $b(v):=\mu(B_G(v,1))$ for the compatible mass at $v$. The potential of the pair
$(G,\mu)$ at exponent $\alpha$ is
\[
 \eta_{G,\alpha}(\mu)=\sum_{v\colon\mu(v)>0}\mu(v)\,\frac{1-b(v)}{b(v)^{\alpha}} .
\]
It is the one-site potential of \cref{def:potential} for the compatibility relation of
\cref{def:compat}.
\end{definition}

\section{The contraction estimate}

\begin{lemma}[$2\times2$ transversal]
\label{thm:transversal}
\lean{GraphMatching.transversal}
\leanok
Mark each cell of a $2\times2$ array good or bad. If no row and no column is entirely bad,
then the array has a good transversal: a good main diagonal or a good antidiagonal.
\end{lemma}

\begin{proof}
\leanok
If the cell $g_{00}$ is good, then either $g_{11}$ is good, giving the good diagonal, or
$g_{11}$ is bad, in which case row $1$ not all bad forces $g_{10}$ good and column $1$ not
all bad forces $g_{01}$ good, giving the good antidiagonal. If $g_{00}$ is bad, then row
$0$ not all bad forces $g_{01}$ good and column $0$ not all bad forces $g_{10}$ good, again
a good antidiagonal.
\end{proof}

\begin{lemma}[Uniform contraction]
\label{thm:contraction}
\lean{GraphMatching.Phi_square_leA}
\leanok
\uses{def:square, thm:transversal}
Let $\alpha\geq1$. For every countable $(\cX,\mu,R)$ with $R$ symmetric and reflexive,
\[
 \Phi(R^\square,\mu^2)\leq A_\alpha\,\Phi(R,\mu)+B_\alpha\,\Phi(R,\mu)^2 .
\]
In particular, if $A_\alpha+B_\alpha\Phi(R,\mu)\leq\rho$ for some $\rho<1$, then
$\Phi(R^\square,\mu^2)\leq\rho\,\Phi(R,\mu)$.
\end{lemma}

\begin{proof}
\leanok
\uses{thm:transversal}
Fix a source pair and arrange its four compatibility tests against a target pair in a
$2\times2$ array. By \cref{thm:transversal}, failure of $R^\square$ forces one row or one
column to be entirely bad. A union bound therefore controls the square bad degree by the two
one-site bad degrees. Applying the pointwise estimates for $\phi_\alpha$ separates this bound
into a linear term and a quadratic remainder. Averaging over the source pair gives the stated
constants $A_\alpha$ and $B_\alpha$.
\end{proof}

\begin{lemma}[Product of relations]
\label{thm:product}
\lean{GraphMatching.PhiA_prodPMF_le}
\leanok
\uses{def:potential}
Let $(\cX_i,\mu_i,R_i)$, $i=1,2$, be countable probability spaces with symmetric reflexive
relations, and let $R_1\otimes R_2$ be the relation on $\cX_1\times\cX_2$ holding when both
coordinates are related. With $\Phi_i:=\Phi(R_i,\mu_i)$ and any $\alpha\geq1$, the
potential of the product is controlled by $\Phi_1$ and $\Phi_2$ at exponent $2\alpha$.
\end{lemma}

\begin{proof}
\leanok
\uses{def:potential}
Fix $(x_1,x_2)$ in the product support and write $q_i$ for the bad degree of $x_i$. Since the
two target coordinates are independent, the good degree of the product point is
$(1-q_1)(1-q_2)$. Its bad degree is therefore
$q_1+(1-q_1)q_2$. The splitting inequality for $\phi_\alpha$ bounds the resulting weight by
one term depending on $q_1$, one depending on $q_2$, and a product term. Summing against
$\mu_1\otimes\mu_2$ yields the claimed control in terms of the two potentials at exponent
$2\alpha$.
\end{proof}

\section{Reduction to the one-site potential}

\begin{theorem}[Reduction]
\label{thm:reduction}
\lean{GraphMatching.leaf_matching_boundA, GraphMatching.full_matching_boundA}
\leanok
\uses{def:matching, thm:contraction, thm:product}
Let $\alpha\geq1$ and let $R_0$ be symmetric and reflexive with one-site potential
$\Phi_0$.
\begin{enumerate}
\item If $\rho:=A_\alpha+B_\alpha\Phi_0<1$, then for leaf labellings
\[
 \bbP(\text{no leaf matching automorphism of }\bbB_h)\leq\Phi_0\,\rho^{\,h}.
\]
\item Let $\mathfrak K>0$ satisfy
$\bigl(1+2\alpha\Phi_0\bigr)\bigl(A_\alpha\mathfrak K+B_\alpha\mathfrak K^2\Phi_0\bigr)
\leq\mathfrak K-1$. Then the full matching probability is bounded below uniformly in the
height $h$.
\end{enumerate}
\end{theorem}

\begin{proof}
\leanok
\uses{thm:contraction, thm:product}
We proceed by induction on the height. For leaf labellings, adjoining one level replaces the
matching relation by its symmetrised square. Thus \cref{thm:contraction} multiplies the
potential by at most $\rho$ at each step, which gives the first bound.

For full labellings, the new root contributes one relation factor and the two independent child
subtrees contribute a symmetrised-square factor. The product estimate and the displayed
hypothesis on $\mathfrak K$ preserve the required uniform bound under this recursion. Finally,
the finite matching sets form a finitely branching, restriction-compatible tree. K\"onig's
lemma produces a compatible automorphism on all levels, while the measure argument preserves
the probability bound. In Lean these steps are
\texttt{Konig.infinite\_matching\_of\_forall\_level} and the accompanying results in
\texttt{Measure}.
\end{proof}

\section{The countable label graph}

\begin{theorem}[Matching for a countable label graph]
\label{thm:matching}
\lean{GraphMatching.graph_leaf_matching_bound, GraphMatching.graph_full_matching_bound}
\leanok
\uses{def:etaG, thm:reduction}
Let $G$ be a countable graph and $\mu$ a law on its vertices.
\begin{enumerate}
\item If $\eta_{G,5/2}(\mu)\leq\tfrac1{256}$, then for leaf labellings
\[
 \bbP(\text{no automorphism of }\bbB_h\text{ matches every leaf})
 \leq\eta_{G,5/2}(\mu)\left(\frac{253}{256}\right)^{\!h},
\]
which tends to zero as $h\to\infty$.
\item If $\eta_{G,5/2}(\mu)\leq10^{-4}$, then for full labellings, uniformly in the height,
\[
 \bbP(\text{some automorphism of }\bbB_h\text{ matches every vertex})
 \geq1-16\,\eta_{G,5/2}(\mu)>0 .
\]
\end{enumerate}
\end{theorem}

\begin{proof}
\leanok
\uses{thm:reduction, def:etaG}
By \cref{def:etaG}, the graph potential is the one-site potential of the compatibility
relation. We may therefore apply \cref{thm:reduction} with
$\Phi_0=\eta_{G,\alpha}(\mu)$. At $\alpha=5/2$, the explicit contraction constants give
$\rho=253/256$ whenever $\Phi_0\leq1/256$. For the full matching recursion,
$\mathfrak K=16$ satisfies the required inequality whenever $\Phi_0\leq10^{-4}$. Substitution
in the two clauses of \cref{thm:reduction} gives the stated bounds.
\end{proof}

\begin{theorem}[The infinite tree]
\label{thm:infinite}
\lean{GraphMatching.exists_infinite_tree_matching_graphAut}
\leanok
Under the hypothesis of the second part of \cref{thm:matching}, the same lower bound holds
for the existence of a single automorphism matching every vertex of two independent
infinite labelled rooted binary trees.
\end{theorem}

\begin{proof}
\leanok
\uses{thm:matching, thm:reduction}
For each height, consider the finite set of matching automorphisms. Restriction from height
$h+1$ to height $h$ preserves the matching property, and \cref{thm:matching} gives the same
lower probability bound at every height. K\"onig's lemma applied to these restriction maps
produces a compatible family of finite automorphisms, hence one automorphism of the infinite
binary tree. Continuity from above transfers the finite-height probability bound to this
intersection event.
\end{proof}

\begin{corollary}[Local-dominance condition]
\label{thm:dominance}
\lean{GraphMatching.etaGA_le_localDominance}
\leanok
\uses{def:etaG}
Fix $0\in V$ with $\mu(0)>0$ and put
\[
 \varepsilon_0:=1-\mu(0),
 \qquad
 \tau:=\mu\bigl(V\setminus B_G(0,1)\bigr),
 \qquad
 \cD_{G,\alpha}(\mu,0):=
 \frac{\tau}{(1-\varepsilon_0)^{\alpha}}
 +\sum_{\substack{v\neq0\\ \mu(v)>0}}\frac{\mu(v)}{\mu(B_G(v,1))^{\alpha}} .
\]
Then $\eta_{G,\alpha}(\mu)\leq\cD_{G,\alpha}(\mu,0)$.
\end{corollary}

\begin{proof}
\leanok
Split the defining sum for $\eta_{G,\alpha}$ into the term at $0$ and the terms with
$v\neq0$. At the distinguished vertex, the incompatible mass is $\tau$ and the compatible
mass is at least $\mu(0)=1-\varepsilon_0$. Its contribution is therefore at most the first
term in $\cD_{G,\alpha}(\mu,0)$. For $v\neq0$, discard the numerator $1-b(v)\leq1$. Summing
these bounds gives the result.
\end{proof}

\section{The integer alphabet}

Specialise to the path graph $\mathsf P$ on the integers, whose potential we write
$\eta_{\mathsf P,\alpha}$.

\begin{corollary}[Double-exponential tail]
\label{thm:double-exp}
\lean{GraphMatching.double_exp_matching}
\leanok
Let $\alpha=5/2$. For $D\geq5$ put $p_j=e^{-D^j}$ for $j\geq1$ and
$p_0=1-\sum_{j\geq1}e^{-D^j}$. Then $p_0>\tfrac12$ and
\[
 \eta_{\mathsf P,5/2}(p)\leq4\,e^{-D(D-5/2)}<10^{-4} .
\]
Consequently the full matching probability is at least
$1-16\eta_{\mathsf P,5/2}(p)\geq1-64\,e^{-D(D-5/2)}>0.9996$, uniformly in the height and
for the infinite tree.
\end{corollary}

\begin{proof}
\leanok
\uses{thm:matching, thm:dominance, thm:infinite}
For $j\geq1$, the ratio of consecutive tail masses satisfies
$p_{j+1}/p_j\leq e^{-20}<1/2$. Hence the sum of the positive-index masses is dominated by a
geometric series, which gives $p_0>1/2$. The same estimate controls the mass outside the unit
ball about $0$, while the neighbouring masses give the denominators required in
\cref{thm:dominance}. Substitution yields
$\eta_{\mathsf P,5/2}(p)\leq4e^{-D(D-5/2)}$. This is below $10^{-4}$ for $D\geq5$, so the
full and infinite-tree conclusions follow from \cref{thm:matching} and \cref{thm:infinite}.
\end{proof}
